\documentclass[12pt, a4paper]{article}
\usepackage{iftex}
\ifPDFTeX
  \usepackage{cmap}
  \usepackage[T2A]{fontenc}
  \usepackage[utf8]{inputenc}
  \usepackage{tempora}
\else
  \usepackage{fontspec}
  \IfFontExistsTF{Times New Roman}
    {\setmainfont{Times New Roman}}
    {\setmainfont{CMU Serif}}
  \IfFontExistsTF{Consolas}
    {\setmonofont{Consolas}[Scale=MatchLowercase]}
    {\IfFontExistsTF{Courier New}
      {\setmonofont{Courier New}}
      {\setmonofont{CMU Typewriter Text}}}
\fi
\usepackage[english,russian]{babel}

\usepackage[a4paper, margin=2.54cm]{geometry}
\setlength{\parindent}{0pt}
\setlength{\parskip}{4pt plus 1pt}
\frenchspacing
\setlength{\emergencystretch}{3em}
\raggedbottom

\usepackage{titlesec}
\setcounter{secnumdepth}{0}
\titleformat{\section}{\large\bfseries}{}{0pt}{}
\titleformat{\subsection}{\normalsize\bfseries}{}{0pt}{}
\titlespacing*{\section}{0pt}{14pt plus 4pt minus 2pt}{4pt plus 2pt}
\titlespacing*{\subsection}{0pt}{10pt plus 2pt minus 2pt}{3pt plus 1pt}

\usepackage[table]{xcolor}
\usepackage{graphicx}
\usepackage{tabularx}
\usepackage{float}
\usepackage{caption}
\DeclareCaptionLabelSeparator{spacedemdash}{~--- }
\captionsetup{labelsep=spacedemdash, font=small, singlelinecheck=false, justification=raggedright}
\captionsetup[table]{position=top, skip=4pt}
\DeclareCaptionType[fileext=lol]{listing}[Листинг][Листинги]
\captionsetup[listing]{position=top, skip=4pt, hypcap=false}
\addto\captionsrussian{\renewcommand{\figurename}{Рисунок}\renewcommand{\tablename}{Таблица}}
\newcommand{\thead}[1]{\multicolumn{1}{|c|}{\cellcolor{gray!20}\textbf{#1}}}
\newcolumntype{L}{>{\raggedright\arraybackslash}X}
\newcolumntype{P}[1]{>{\raggedright\arraybackslash}p{#1}}

\usepackage{fvextra}
\usepackage{needspace}
\fvset{breaklines=true, breakanywhere=true,
       frame=leftline, framerule=1.5pt, framesep=6pt,
       rulecolor=\color{gray!60}, xleftmargin=4pt}
\DefineVerbatimEnvironment{code}{Verbatim}{fontsize=\small}
\DefineVerbatimEnvironment{terminal}{Verbatim}{fontsize=\footnotesize}
\newcommand{\inputterminal}[2][]{\VerbatimInput[fontsize=\footnotesize,#1]{#2}}
\ifPDFTeX
  \usepackage{listingsutf8}
  \lstset{inputencoding=utf8/cp1251, basicstyle=\fontencoding{T2A}\selectfont\ttfamily\small,
          breaklines=true, breakatwhitespace=false, columns=fullflexible,
          keepspaces=true, showstringspaces=false, frame=leftline,
          framerule=1.5pt, framesep=6pt, xleftmargin=4pt,
          rulecolor=\color{gray!60}}
  \newcommand{\inputcode}[2][]{\lstinputlisting[#1]{#2}}
  \newcommand{\inputappendixcode}[2][]{\lstinputlisting[basicstyle=\linespread{0.90}\fontencoding{T2A}\selectfont\ttfamily\footnotesize,#1]{#2}}
\else
  \newcommand{\inputcode}[2][]{\VerbatimInput[fontsize=\small,#1]{#2}}
  \newcommand{\inputappendixcode}[2][]{\VerbatimInput[fontsize=\footnotesize,baselinestretch=0.90,#1]{#2}}
\fi
\newenvironment{listingblock}[1]
  {\Needspace{8\baselineskip}\par\medskip\captionof{listing}{#1}\nopagebreak}
  {\par}

\newcounter{appendixno}
\newcommand{\appendixsection}[2]{\clearpage\section{Приложение #1. #2}\stepcounter{appendixno}\setcounter{listing}{0}\renewcommand{\thelisting}{#1.\arabic{listing}}\renewcommand{\theHlisting}{app\arabic{appendixno}.\arabic{listing}}}

\usepackage[hidelinks]{hyperref}
\usepackage{xurl}

\newcommand{\hint}[1]{\leavevmode{\color{gray}\itshape #1}}
\newcommand{\placeholder}[2][4.5cm]{\fbox{\parbox[c][#1][c]{0.95\linewidth}{\centering\hint{#2}}}}


\newcommand{\ReportLab}{5}
\newcommand{\ReportStudent}{Новак Дмитрий Евгеньевич}
\newcommand{\ReportGroup}{М8О-201БВ-25}
\newcommand{\ReportTeacher}{Миронов Евгений Сергеевич}
\newcommand{\ReportDate}{26.11.2026}
\newcommand{\ReportYear}{2026}
\newcommand{\ReportCourse}{Операционные системы и архитектура компьютера}

\begin{document}
% Титульный лист. Данные берутся из команд \Report... в report.tex.
\begin{titlepage}
  \centering\large

  МОСКОВСКИЙ АВИАЦИОННЫЙ ИНСТИТУТ\\
  (НАЦИОНАЛЬНЫЙ ИССЛЕДОВАТЕЛЬСКИЙ УНИВЕРСИТЕТ)\par
  \vspace{18pt}
  Институт №8 «Компьютерные науки и прикладная математика»\\
  Кафедра 806 «Вычислительная математика и программирование»\par

  \vspace{\stretch{1.2}}

  {\bfseries Лабораторная работа №\,\ReportLab{} по курсу\\
  «\ReportCourse»\par}

  \vspace{\stretch{1.6}}

  \begin{flushright}
    \begin{tabular}{r@{\hspace{0.6em}}l}
      Выполнил:      & \ReportStudent \\[2pt]
      Группа:        & \ReportGroup   \\[2pt]
      Преподаватель: & \ReportTeacher \\[2pt]
      Дата:          & \ReportDate    \\[2pt]
      Оценка:        & \rule[-2pt]{4.2cm}{0.4pt} \\
    \end{tabular}
  \end{flushright}

  \vspace{\stretch{1}}

  Москва, \ReportYear
\end{titlepage}
\setcounter{page}{2}


\section{Условие}

\textbf{Цель работы:} получить практические навыки создания динамических библиотек и программ, использующих их функции.

\textbf{Задание:} создать две динамические библиотеки с одинаковыми контрактами и разными реализациями. Написать две программы: первая связывается с одной библиотекой при сборке, вторая загружает библиотеки во время выполнения и переключается между ними по команде пользователя. Сравнить два способа использования библиотек.

\textbf{Вариант 20:} реализовать подсчёт простых чисел на отрезке натуральных чисел $[A, B]$ и перевод целого числа из десятичной системы счисления. Первая библиотека проверяет делители перебором и переводит число в двоичную систему; вторая использует решето Эратосфена и переводит число в троичную систему.

\section{Метод решения}

Обе библиотеки предоставляют функции \texttt{PrimeCount()} и \texttt{translation()}. Первая принимает две границы типа \texttt{int}, вторая --- число типа \texttt{long}. Первая возвращает количество простых чисел на отрезке включительно или отрицательное значение при ошибке. Вторая выделяет память для строки с результатом; вызывающая программа освобождает её с помощью \texttt{free()}.

В первой реализации каждое число проверяется перебором возможных делителей от 2 до числа минус один. Если найден делитель, число не является простым. Во второй реализации создаётся массив признаков составных чисел: для каждого непомеченного числа отмечаются его кратные, после чего подсчитываются непомеченные числа на заданном отрезке. Решето использует дополнительную память, зато не проверяет делители каждого числа отдельно.

При переводе числа последовательно находятся остатки от деления на основание системы счисления: 2 или 3. Полученные цифры записываются в обратном порядке. Для отрицательного числа перед результатом ставится знак минус; ноль обрабатывается как отдельная цифра.

Программа \texttt{lab5\_linked} связывается с первой библиотекой при сборке. Программа \texttt{lab5\_dynamic} загружает первую библиотеку при запуске через \texttt{dlopen()}, получает адреса функций через \texttt{dlsym()} и сохраняет их в структуре \texttt{Library}. По команде \texttt{0} загружается другая библиотека, предыдущая закрывается через \texttt{dlclose()}, а вызовы направляются на новые функции. Перед выходом программа закрывает текущую библиотеку.

При связывании зависимость от первой библиотеки задаётся на этапе сборки, а её загрузку при запуске обеспечивает динамический загрузчик. При явной загрузке программа сама выбирает библиотеку и может переключить реализацию во время работы, но должна проверять результат загрузки и поиска функций. В обоих случаях используются динамические библиотеки; функции \texttt{dlopen()} и \texttt{dlsym()} не являются отдельными системными вызовами.

\section{Описание программы}

\begin{table}[H]
  \caption{Файлы программной части}
  \begin{tabularx}{\linewidth}{|P{4.2cm}|L|}
    \hline
    \thead{Файл} & \thead{Назначение} \\ \hline
    \path{src/lib1.c} & Перебор делителей и перевод числа в двоичную систему. \\ \hline
    \path{src/lib2.c} & Решето Эратосфена и перевод числа в троичную систему. \\ \hline
    \path{src/program1.c} & Программа со связыванием первой библиотеки при сборке. \\ \hline
    \path{src/program2.c} & Загрузка библиотек, поиск функций и переключение реализаций. \\ \hline
    \texttt{CMakeLists.txt} & Сборка библиотек и двух исполняемых программ. \\ \hline
    \path{report/trace/*.txt} & Трассировки двух запусков в Ubuntu ARM64. \\ \hline
  \end{tabularx}
\end{table}

Обе программы читают команды построчно с помощью \texttt{getline()}. Функция \texttt{read\_long()} использует \texttt{strtol()} для преобразования и проверки чисел. Проверяются команда, число аргументов и диапазон границ; при ошибке выводится сообщение и программа завершается.

\begin{table}[H]
  \caption{Команды пользователя}
  \begin{tabularx}{\linewidth}{|P{2.8cm}|L|}
    \hline
    \thead{Команда} & \thead{Действие} \\ \hline
    \texttt{1 A B} & Подсчитать простые числа на отрезке $[A, B]$. \\ \hline
    \texttt{2 x} & Перевести число в систему счисления текущей реализации. \\ \hline
    \texttt{0} & Переключить библиотеку во второй программе. \\ \hline
  \end{tabularx}
\end{table}

\section{Инструкция по запуску}

Сборка из корня репозитория:
\begin{terminal}
$ cmake -S lab_05 -B build/lab_05
$ cmake --build build/lab_05
$ cd build/lab_05
\end{terminal}

В исходном файле \texttt{program2.c} указаны имена библиотек для macOS с расширением \texttt{.dylib}. Для запуска второй программы в Linux перед сборкой необходимо заменить оба расширения на \texttt{.so}. Библиотеки загружаются по относительным путям, поэтому программа запускается из каталога сборки.

Первая программа использует первую реализацию:
\begin{terminal}
$ printf '1 1 10\n2 10\n' | ./lab5_linked
4
1010
\end{terminal}

Во второй программе команда \texttt{0} переключает библиотеку:
\begin{terminal}
$ printf '1 1 10\n2 10\n0\n1 1 10\n2 10\n' | ./lab5_dynamic
4
1010
4
101
\end{terminal}

На отрезке $[1, 10]$ четыре простых числа: 2, 3, 5 и 7. Число 10 записывается как $1010_2$ и $101_3$.

\section{Демонстрация работы strace}

Трассировки получены на Ubuntu ARM64 после замены расширений библиотек на \texttt{.so}. Команды выполнялись из каталога \texttt{build/lab\_05}:
\begin{terminal}
$ printf '1 1 10\n2 10\n' | strace -o linked.txt -s 256 ./lab5_linked
$ printf '1 1 10\n2 10\n0\n1 1 10\n2 10\n' | strace -o dynamic.txt -s 256 ./lab5_dynamic
\end{terminal}

В первой трассировке загрузчик открывает \texttt{liblab5\_first.so} до чтения команд. Во второй программе сначала открывается первая библиотека, после ввода и первых двух результатов открывается \texttt{liblab5\_second.so}, а отображение первой освобождается через \texttt{munmap()}. Затем выводятся результаты второй реализации. Обе программы завершаются с кодом 0.

\begin{listingblock}{Выбранные строки первой трассировки}
  \inputterminal[firstline=46,lastline=52]{trace/linked.txt}
\end{listingblock}

\begin{listingblock}{Выбранные строки второй трассировки}
  \inputterminal[firstline=46,lastline=50]{trace/dynamic.txt}
  \inputterminal[firstline=62,lastline=68]{trace/dynamic.txt}
\end{listingblock}

Вызовы \texttt{read()} и \texttt{write()} показывают ввод и результаты; \texttt{openat()} и \texttt{mmap()} связаны с открытием и отображением библиотек. По одним строкам трассировки без стека вызовов нельзя однозначно определить конкретную строку исходного кода. Полные трассировки приведены в приложении Б.

\Needspace{8\baselineskip}
\section{Выводы}

Таким образом, данная лабораторная работа никак кардинально не изменила мою жизнь: я научился создавать и подключать динамические библиотеки. Теперь у меня есть две реализации одних и тех же функций, которые умеют договариваться с одной программой, не зная друг о друге.

\appendixsection{А}{Исходный код программы}
\begin{listingblock}{Файл lib1.c}
  \inputappendixcode{../src/lib1.c}
\end{listingblock}
\begin{listingblock}{Файл lib2.c}
  \inputappendixcode{../src/lib2.c}
\end{listingblock}
\begin{listingblock}{Файл program1.c}
  \inputappendixcode{../src/program1.c}
\end{listingblock}
\begin{listingblock}{Файл program2.c}
  \inputappendixcode{../src/program2.c}
\end{listingblock}
\begin{listingblock}{Файл CMakeLists.txt}
  \inputappendixcode{../CMakeLists.txt}
\end{listingblock}

\appendixsection{Б}{Полный вывод strace}
\begin{listingblock}{Трассировка lab5\_linked}
  \inputterminal[baselinestretch=0.95]{trace/linked.txt}
\end{listingblock}
\begin{listingblock}{Трассировка lab5\_dynamic}
  \inputterminal[baselinestretch=0.95]{trace/dynamic.txt}
\end{listingblock}
\end{document}
